\chapter{Ergebnisse}
\label{chap:results}

Zur Orientierung gilt: Die Pipeline-Status \textit{Attempted}, \textit{Implemented\textsubscript{static}} und \textit{Verified} beschreiben unterschiedliche Evidenzebenen desselben Arbeitsprozesses und sind keine Qualitäts- oder Kompetenzbewertung. Insbesondere gilt: \textit{Verified} ist nicht gleich \enquote{besserer Entwickler}. \textit{Implemented\textsubscript{static}} meint strukturelle Artefakte in der Codebasis; \textit{Verified} ist strikt testgebunden.

\section{Datengrundlage und Toolkontext}
\label{sec:results_context}

Dieses Unterkapitel beschreibt die Datengrundlage der Untersuchung sowie den technischen Kontext der durchgeführten \textit{VibeCoding-Challenge}. Die Analyse basiert auf einem Datensatz von $N=18$ Probanden, die im Rahmen einer kontrollierten Coding-Challenge Aufgaben zur Entwicklung einer Todo-Applikation bearbeiteten.

\subsection{Studiensetting und Rahmenbedingungen}
Die Studie wurde unter einheitlichen Rahmenbedingungen durchgeführt, um die Vergleichbarkeit der Ergebnisse zu gewährleisten:
\begin{itemize}
    \item \textbf{Zeitlimit:} Alle Probanden hatten ein striktes Zeitlimit von 60 Minuten für die Bearbeitung der Aufgaben.
    \item \textbf{Tooling:} Als Entwicklungsumgebung diente Visual Studio Code mit der GitHub Copilot Extension.
    \item \textbf{KI-Modell:} Als zugrundeliegendes Large Language Model (LLM) wurde \textbf{Claude 3.5 Sonnet} verwendet. Die Interaktion erfolgte über das Chat-Interface von GitHub Copilot, wobei das Modell explizit als Backend konfiguriert war.
    \item \textbf{Aufgabenstellung:} Die Probanden erhielten eine Liste von 8 Aufgaben (T1--T8) mit steigender Komplexität, beginnend bei Backend-Persistenz bis hin zu komplexen Frontend-Features und Export-Funktionen.
\end{itemize}

\subsection{Datenquellen und Artefakte}
Die empirische Auswertung stützt sich auf drei primäre Datenquellen, die eine Triangulation der Ergebnisse ermöglichen:

\begin{enumerate}
    \item \textbf{Code-Artefakte (Git-Repositories):}
    Die finalen Implementierungen der Probanden liegen als separate Git-Branches vor. Diese Branches enthalten den vollständigen Quellcode zum Zeitpunkt des Zeitablaufs.
    \begin{itemize}
        \item Pfad: Remote-Branches mit dem Suffix \texttt{*-vibe} (z.B. \texttt{origin/august-martin-vibe}).
    \end{itemize}

    \item \textbf{Chat-Logs:}
    Die vollständigen Interaktionsprotokolle zwischen Proband und KI-Assistent wurden exportiert. Diese Protokolle enthalten alle Prompts der Nutzer sowie die Antworten und Code-Generierungen der KI.
    \begin{itemize}
        \item Pfad: \texttt{thesis/chats/} (Format: \texttt{.docx}).
    \end{itemize}

    \item \textbf{Survey-Daten:}
    Im Anschluss an die Challenge füllten die Probanden einen Fragebogen zur Selbsteinschätzung und subjektiven Wahrnehmung der KI-Interaktion aus.
    \begin{itemize}
        \item Pfad: \texttt{thesis/surveys/Probanden-Einstufungsformular (Responses).xlsx}.
    \end{itemize}
\end{enumerate}

\subsection{Dateninventar und Vollständigkeit}
Tabelle \ref{tab:data_inventory} zeigt das Inventar der verfügbaren Datenartefakte. Insgesamt liegen für 18 Probanden Daten vor.

\begin{table}[h]
    \centering
    \begin{tabular}{|l|c|l|}
    \hline
    \textbf{Artefakt-Typ} & \textbf{Anzahl} & \textbf{Anmerkung} \\
    \hline
    Git-Branches (\texttt{*-vibe}) & 16 & 2 Probanden arbeiteten lokal/ohne Push* \\
    Chat-Logs & 18 & Vollständig für alle Teilnehmer \\
    Survey-Einträge & 18 & Vollständig ($N=18$) \\
    \hline
    \end{tabular}
    \caption{Inventar der Datenartefakte. (*Für die fehlenden Remote-Branches liegen lokale Stände oder Chat-Rekonstruktionen vor, die für die Analyse genutzt wurden.)}
    \label{tab:data_inventory}
\end{table}

\subsection{Zuordnung und Zusammenführung der Daten}
Die Zusammenführung der Datenquellen erfolgte über eine Mapping-Tabelle, die die Klarnamen aus dem Survey und den Chat-Dateinamen den entsprechenden Git-Branches zuordnet.
\begin{itemize}
    \item \textbf{Survey $\leftrightarrow$ Branch:} Das Survey-Feld "Branch-Name" sowie der Name des Probanden dienten als Schlüssel.
    \item \textbf{Chat $\leftrightarrow$ Branch:} Die Dateinamen der Chat-Logs (z.B. \texttt{Max\_Helling.docx}) wurden normalisiert und den entsprechenden Branches (z.B. \texttt{helling-max-vibe}) zugeordnet.
\end{itemize}
Die Survey-Daten enthalten zudem eine Einteilung in Kompetenzgruppen (No-Code, Low-Code, Mid-Code, High-Code), die in der späteren Analyse als vergleichendes Merkmal herangezogen wird.

\subsection{Reproduzierbarkeit und Auswertungswerkzeuge}
Um die Ergebnisse objektiv und reproduzierbar zu gestalten, kommen automatisierte Analysewerkzeuge zum Einsatz, deren Resultate im Ordner \texttt{evaluation/} persistiert werden:

\begin{itemize}
    \item \textbf{Statische Code-Analyse:} Einsatz von \texttt{ESLint} zur Überprüfung der Code-Qualität und Einhaltung von Standards.
    \item \textbf{Automatisierte Tests:} Ausführung der vorhandenen Test-Suite (basierend auf \texttt{Jest/Vitest}), um die funktionale Korrektheit der Implementierungen (Task Completion) zu verifizieren.
    \item \textbf{Chat-Parsing:} Algorithmische Auswertung der Chat-Logs zur Bestimmung von "Attempted"-Tasks und Identifikation von Prompting-Mustern.
\end{itemize}

\subsection{Abgrenzung zu Kapitel 7}
Das vorliegende Kapitel 6 beschränkt sich auf die deskriptive Darstellung der Ergebnisse und die objektive Auswertung der Daten. Eine Interpretation dieser Ergebnisse, die Diskussion von Ursache-Wirkungs-Beziehungen sowie die Beantwortung der Forschungsfragen erfolgen im anschließenden Kapitel 7 (Diskussion).

\section{Task-Pipeline-Analyse}
\label{sec:results_pipeline}

In diesem Abschnitt wird der Fortschritt der Probanden entlang der Aufgaben-Pipeline (T1--T5) dargestellt. Berichtete werden Häufigkeiten für Intention (Chat), Umsetzung (Code-Artefakte) und den Statusausweis \textit{Execution observed} (Integrationsläufe).

\subsection{Methodische Definitionen}
Für die Auswertung werden drei Statuskategorien verwendet:
\begin{itemize}
    \item  \textbf{Attempted (Versucht):} Ein Task gilt als  \textit{Attempted}, wenn er im Chat-Verlauf explizit als Ziel, Anforderung oder nächster Arbeitsschritt formuliert wurde. Die bloße Existenz von Dateien ohne entsprechenden Kontext im Chat genügt nicht.
    \item  \textbf{Implemented\textsubscript{static} (Umgesetzt, Artefakt/Heuristik):}\\
    Ein Task gilt als  \textit{Implemented\textsubscript{static}}, wenn er in den Auswertungsartefakten als implementiert markiert ist (heuristische, statische Indikatoren; z.\,B.\ Feld \texttt{implemented}\\
    in \path{evaluation/task_pipeline_v3.json}\\
    bzw.\ Implementationsklassen in \mbox{\path{evaluation/deep_code_analysis.json}}).
    \item  \textbf{Execution observed (Statusausweis):} Ein Task gilt als  \textit{Execution observed}, wenn die zugehörigen automatisierten Integrationsläufe erfolgreich durchlaufen wurden (Feld \texttt{verified} in \path{evaluation/task_pipeline_v3.json}).
\end{itemize}

\noindent \textbf{Abgeleitete Evidenzstufe.} In Kapitel~6 wird eine abgeleitete Evidenzstufe\\
\BeginAccSupp{unicode,method=pdfstringdef,ActualText={extitImplemented* verwendet:}}
\csname textit\endcsname{Implemented$^\ast$} verwendet:
\EndAccSupp{}
\[
\begin{aligned}
    \csname textit\endcsname{Implemented$^\ast$} &:=  \textit{Attempted} \wedge (\textit{Implemented\textsubscript{static}} \vee  \textit{Execution observed}).
\end{aligned}
\]
Damit gilt per Konstruktion:  \textit{Execution observed} $\subseteq$  \textit{Implemented$^\ast$}.  \textit{Execution observed} ist jedoch nicht notwendigerweise ein Subset von  \textit{Implemented\textsubscript{static}}:  \textit{Implemented\textsubscript{static}} basiert auf heuristischen, strukturellen Erkennungsregeln; unter den gegebenen Rahmenbedingungen ist es möglich, dass funktionsfähige Implementationen nicht als  \textit{Implemented\textsubscript{static}} erkannt werden (z.\,B. abweichende Benennungen, alternative Architekturpfade oder Implementationen außerhalb der erwarteten Muster).  \textit{Execution observed} ist ein methodengebundenes Diagnose-Signal;  \textit{Implemented\textsubscript{static}} ist ein Artefakt-Signal.

\noindent \textbf{Hinweis zur Ausführungsmethodik.} Für die Pipeline-Auswertung wurden die automatisierten Integrationsläufe (insbesondere T2, sowie in Teilen T3--T5) als API-Integrationsläufe (\textit{Supertest} gegen Express) gestaltet. \textit{Execution observed} bedeutet in dieser Studie explizit: „durchläuft die aktuellen Integrationsläufe“ \textendash{} nicht „fachlich vollständig korrekt gelöst“.

\subsection{Pipeline-Ergebnisse (T1--T5)}
\begin{table}[ht]
\centering
\begingroup
\newcommand{\NoExtract}[1]{\BeginAccSupp{method=escape,ActualText={}}#1\EndAccSupp{}}
\newcommand{\ExtractOnly}[1]{\smash{\rlap{\BeginAccSupp{unicode,method=pdfstringdef,ActualText={#1}}\textcolor{white}{.}\EndAccSupp{}}}}
\def\TaskPipelineTableActualText{Task^^JAttempted (𝑛) Implemented* (𝑛) Execution observed (𝑛)^^JT1: Persistenz^^J18 (100\%)^^J17 (94\%)^^J10 (56\%)^^JT2: Auth^^J18 (100\%)^^J17 (94\%)^^J4 (22\%)^^JT3: Kategorien^^J15 (83\%)^^J13 (72\%)^^J13 (72\%)^^JT4: Attachments^^J13 (72\%)^^J8 (44\%)^^J0 (0\%)^^JT5: Kalender^^J13 (72\%)^^J0 (0\%)^^J0 (0\%)^^J^^JTabelle 6.2: Task-Pipeline: Von der Intention zur testgebundenen Bestätigung (𝑁 = 18,^^J𝑛 = Anzahl Branches).^^JextitAttempted stammt aus Chat-Signalen. Implemented* ist definiert als: Attempted und^^J(Implementedstatic oder Execution observed). Execution observed stammt aus V3-Integrationsläufen^^J(V3 bezeichnet die für diese Studie definierte Integrationslauf-Konfiguration).}
\NoExtract{\begin{tabular}{|l|c|c|c|}
\hline
    \csname textbf\endcsname{Task} &  \textbf{Attempted ($n$)} &  \textbf{Implemented$^\ast$ ($n$)} &  \textbf{Execution observed ($n$)} \\
\hline
T1: Persistenz & 18 (100\%) & 17 (94\%) & 10 (56\%) \\
T2: Auth & 18 (100\%) & 17 (94\%) & 4 (22\%) \\
T3: Kategorien & 15 (83\%) & 13 (72\%) & 13 (72\%) \\
T4: Attachments & 13 (72\%) & 8 (44\%) & 0 (0\%) \\
T5: Kalender & 13 (72\%) & 0 (0\%) & 0 (0\%) \\
\hline
\end{tabular}}
\ExtractOnly{\TaskPipelineTableActualText}
\NoExtract{\caption{Task-Pipeline: Von der Intention zur testgebundenen Bestätigung ($N = 18$, $n$ = Anzahl Branches). \textit{Attempted} stammt aus Chat-Signalen. \textit{Implemented$^\ast$} ist definiert als: \textit{Attempted} und (\textit{Implemented\textsubscript{static}} oder \textit{Execution observed}). \textit{Execution observed} stammt aus V3-Integrationsläufen (V3 bezeichnet die für diese Studie definierte Integrationslauf-Konfiguration).}}
\label{tab:task_pipeline_v2}
\endgroup
\end{table}

\noindent \textbf{Messung.} Berichtete werden je Task (T1--T5) die Häufigkeiten der Pipeline-Status \textit{Attempted} (Chat-Signal), \textit{Implemented$^\ast$} (kombinierte Evidenzstufe) und \textit{Execution observed} (methodengebundenes Signal) über $N=18$ Branches.
\noindent \textbf{Hauptbefund.} Für T1/T2 liegt \textit{Attempted} jeweils in 18/18 Branches vor, während \textit{Execution observed} für T4/T5 in 0/18 Branches erreicht wird.
\noindent \textbf{RQ-Bezug.} Die Statushäufigkeiten bilden die deskriptive Ergebnisbasis für RQ1 und werden in Kapitel~7 für die Einordnung der Befunde und die Pipeline-Diskussion verwendet.

\begin{figure}[ht]
\centering
\makebox[\linewidth][c]{%
\begin{tikzpicture}
\begin{axis}[
    xbar,
    width=0.92\linewidth,
    height=7.2cm,
    xmin=0,
    xmax=18,
    xtick={0,2,...,18},
    scaled x ticks=false,
    tick label style={/pgf/number format/fixed,/pgf/number format/precision=0},
    xticklabel style={font=\scriptsize},
    xlabel={Anzahl Branches ($N=18$)},
    xlabel style={font=\small},
    symbolic y coords={T1: Persistenz,T2: Auth,T3: Kategorien,T4: Attachments,T5: Kalender},
    ytick=data,
    y dir=reverse,
    bar width=6pt,
    enlarge y limits=0.45,
    trim axis left,
    trim axis right,
    axis lines=left,
    xmajorgrids,
    grid style={black!8, line width=0.15pt},
    tick style={black!50},
    legend style={
        at={(0.5,-0.18)},
        anchor=north,
        legend columns=3,
        font=\small,
        draw=none
    },
    point meta=explicit symbolic,
    nodes near coords,
    nodes near coords style={
        font=\scriptsize,
        anchor=west,
        xshift=3pt,
        text=black!80
    },
    clip=false,
]
    % Attempted (hellgrau)
    \addplot[fill=black!15, draw=black!55, line width=0.2pt, bar shift=9pt]
    coordinates {(18,T1: Persistenz) [18] (18,T2: Auth) [18] (15,T3: Kategorien) [15] (13,T4: Attachments) [13] (13,T5: Kalender) [13]};

    % Implemented* (mittelgrau, visuell dominanter)
    \addplot[fill=black!55, draw=black, line width=0.35pt, bar width=7pt, bar shift=0pt]
    coordinates {(17,T1: Persistenz) [17] (17,T2: Auth) [17] (13,T3: Kategorien) [13] (8,T4: Attachments) [8] (0,T5: Kalender) [{}]};

    % Execution observed (dunkelgrau/schwarz; zusätzlich Schraffur)
    \addplot[fill=black!85, draw=black, line width=0.4pt, pattern=north east lines, pattern color=black, bar shift=-9pt]
    coordinates {(10,T1: Persistenz) [10] (4,T2: Auth) [4] (13,T3: Kategorien) [13] (0,T4: Attachments) [{}] (0,T5: Kalender) [{}]};

    % 0-Werte (ohne Textlabel) optional als Marker am Achsenstart
    \addplot[only marks, mark=o, mark size=2.6pt, draw=black, mark layer=foreground, forget plot, bar shift=0pt]
    coordinates {(0,T5: Kalender)};
    \addplot[only marks, mark=x, mark size=3.0pt, draw=black, mark layer=foreground, forget plot, bar shift=-9pt]
    coordinates {(0,T4: Attachments) (0,T5: Kalender)};

    \legend{Attempted, Implemented$^\ast$, Execution observed}
\end{axis}
\end{tikzpicture}
}%
\caption{Pipeline-Ergebnisse für T1--T5 ($N=18$) als gruppierte Balken:  \textit{Attempted} (Chat-Signal),  \textit{Implemented$^\ast$} (definiert als:  \textit{Attempted} und (\textit{Implemented\textsubscript{static}} oder \textit{Execution observed})) und  \textit{Execution observed} (testgebunden: bestandene V3-Integrationsläufe). V3 bezeichnet die für diese Studie definierte Integrationslauf-Konfiguration; die Werte entsprechen Tabelle~\ref{tab:task_pipeline_v2}.}
\label{fig:task_pipeline_funnel_t1_t5}
\end{figure}

\noindent \textbf{Messung.} Visualisiert werden die in Tabelle~\ref{tab:task_pipeline_v2} berichteten Statushäufigkeiten (\textit{Attempted}, \textit{Implemented$^\ast$}, \textit{Execution observed}) als Balken pro Task.
\noindent \textbf{Hauptbefund.} Über T1--T5 liegen \textit{Implemented$^\ast$} und \textit{Execution observed} unter \textit{Attempted} und erreichen bei T4/T5 den Wert 0.
\noindent \textbf{RQ-Bezug.} Die Abbildung ist eine komprimierte Darstellung der Pipeline-Ergebnisse und ergänzt die Ergebniszusammenfassung zu RQ1 in Kapitel~7.

\subsection{Kurzbefunde pro Task (T1--T5)}
\noindent \textbf{T1 (Persistenz).} T1 ist in 18/18 Branches  \textit{Attempted}, in 17/18 Branches  \textit{Implemented$^\ast$} und in 10/18 Branches  \textit{Execution observed}.

\noindent \textbf{T2 (User Accounts/Auth).} T2 ist in 18/18 Branches  \textit{Attempted}, in 17/18 Branches  \textit{Implemented$^\ast$} und in 4/18 Branches  \textit{Execution observed}.

\noindent \textbf{T3 (Kategorien).} T3 ist in 15/18 Branches  \textit{Attempted}, in 13/18 Branches  \textit{Implemented$^\ast$} und in 13/18 Branches  \textit{Execution observed}. In \path{evaluation/task_pipeline_v3.json} ist T3 dabei in 7/18 Branches als \texttt{implemented} (\textit{Implemented\textsubscript{static}}) markiert.

\noindent \textbf{T4 (Attachments).} T4 ist in 13/18 Branches  \textit{Attempted}, in 8/18 Branches  \textit{Implemented$^\ast$} und in 0/18 Branches  \textit{Execution observed}.

\noindent \textbf{T5 (Kalender/Multi-Day).} T5 ist in 13/18 Branches  \textit{Attempted}, in 0/18 Branches  \textit{Implemented$^\ast$} und in 0/18 Branches  \textit{Execution observed}.

% DATENQUELLEN:
% Chat-Logs (für "Attempted"-Status)
% Code-Analyse (für "Implemented"-Status)
% Test-Reports (für "Execution observed"-Status)
%
% HINWEISE:
% Definition "Attempted": Explizite Erwähnung im Chat
% Definition "Implemented": Code vorhanden
% Definition "Execution observed": Tests bestanden
% Fokus auf die Statushäufigkeiten über die Tasks hinweg

\section{Detaillierte Lösungsmuster pro Task (T1--T5)}
\label{sec:results_strategies}

\BeginAccSupp{unicode,method=pdfstringdef,ActualText={Dieses Unterkapitel beschreibt typische Lösungsmuster, die in den 𝑁 = 18 *-vibeBranches für die Tasks T1–T5 beobachtet wurden. Grundlage sind Repository-Struktur^^Jund Implementationsartefakte (Backend/Frontend, Datenhaltung), Chat-basierte Attempted-Signale sowie Survey-Angaben (aggregiert, ohne Klarnamen). Die Darstellung ist^^Jrein deskriptiv.}}
Dieses Unterkapitel beschreibt typische Lösungsmuster, die in den $N = 18$ *-vibeBranches für die Tasks T1--T5 beobachtet wurden. Grundlage sind Repository-Struktur und Implementationsartefakte (Backend/Frontend, Datenhaltung), Chat-basierte  \textit{Attempted}-Signale sowie Survey-Angaben (aggregiert, ohne Klarnamen). Die Darstellung ist rein deskriptiv.
\EndAccSupp{}

\subsection*{Einordnung der Datenbasis (kurz)}
T1 und T2 wurden im Chat durchgängig adressiert (18/18), T3 in 15/18 sowie T4/T5 in jeweils 13/18 Fällen. Entsprechend nimmt die Abdeckung der Chat-basierten  \textit{Attempted}-Signale von T1/T2 zu T4/T5 ab. In den Surveys werden Zeitdruck und Debugging am häufigsten als Hürden genannt (je 5/18), gefolgt von Aufgabenverständnis (3/18). Die Selbsteinstufung verteilt sich über No-Code (5), Low-Code (4), Mid-Code (5) und High-Coder (4).

\subsection*{T1 -- Persistente Todos}
\begin{enumerate}
  \item  \textbf{Ziel des Tasks (kurz, technisch).} Persistente Speicherung von Todos über Prozess-Neustarts hinweg, inkl. stabilem CRUD-Verhalten (Erstellen, Lesen, Aktualisieren, Löschen) und konsistenter Identifikatoren.

  \item  \textbf{Typische Implementierungsvarianten (mit Häufigkeiten).} Die häufigste Variante war eine SQLite-basierte Persistenz (14/18). Daneben traten MongoDB (1/18) und PostgreSQL (1/18) sowie verbleibende Varianten wie In-Memory/sonstige Ablagen (2/18) auf. In der Code-Struktur wird häufig ein Repository-Pattern (z.B. \texttt{TodoRepository}) mit einer klaren Abgrenzung zwischen Controller/Route und Datenzugriff verwendet.

  \item  \textbf{Beobachtete Stabilitätsmuster (konsistent vs. inkonsistent).} Konsistent sind konsistente ID-Schemata und deterministisches CRUD. Inkonsistent sind v.a. uneinheitliche ID-Felder (\texttt{id} vs. \texttt{\_id}) und inkonsistente Update-Semantik.

  \item  \textbf{Beobachtete Implementationsmuster.} 
  \begin{sloppypar}Häufig treten Repository-/Controller-Boilerplates und generische CRUD-Endpunkte auf. Randfälle (Eingabeprüfung, Fehlercodes, ID-Normalisierung) sind in den Artefakten nicht in allen Branches durchgängig konsistent umgesetzt.\end{sloppypar}

  \item  \textbf{Beobachtete Teilumsetzungen / offene Punkte.} In einzelnen Branches bleiben Randfälle (Eingabeprüfung, Fehlercodes, ID-Normalisierung) sowie Konsistenzanforderungen (einheitliche ID-Felder, deterministische Update-Semantik) unvollständig.
\end{enumerate}

\subsection*{T2 -- Benutzerkonten / Authentifizierung}
\begin{enumerate}
  \item  \textbf{Ziel des Tasks (kurz, technisch).} Nutzerregistrierung und Login mit Zugriffskontrolle auf geschützte Endpunkte (z.B. \texttt{GET /api/todos}) sowie konsistenter Authentifizierungsweitergabe (Token oder Session).

  \item  \textbf{Typische Implementierungsvarianten (mit Häufigkeiten).} JWT ist die häufigste Variante (11/18), gefolgt von Session/Cookie (4/18) und Hybrid-Ansätzen (1/18). Selten sind Sonderfälle (z.B. Hashing ohne klaren Token-Flow: 1/18) bzw. fehlende oder unklare Signale (1/18). Häufige Strukturmuster sind zentrale Auth-Middleware oder Kontrollen direkt im Controller.

  \item  \textbf{Beobachtete Stabilitätsmuster (konsistent vs. inkonsistent).} Konsistent sind eindeutige Token-/Session-Flows und konsistente Route-Guards. Inkonsistent sind inkonsistente Header-/Cookie-Nutzung und divergierende Erwartungen zwischen Frontend und Backend.

  \item  \textbf{Beobachtete Implementationsmuster.}\\
  Boilerplate-Flows (Register, Login, Protect) treten häufig auf. In den Artefakten finden sich sowohl Varianten mit Rollen-/Claims-Strukturen als auch minimal gehaltene Flows (z.B. ohne durchgängige Eingabeprüfung/Expiry). Häufig treten Kombinationen aus JWT und \texttt{bcrypt} auf.

  \item  \textbf{Beobachtete Teilumsetzungen / offene Punkte.} In einzelnen Branches sind Header-/Cookie-Nutzung und Route-Guards inkonsistent oder zwischen Frontend und Backend nicht durchgängig abgestimmt.
\end{enumerate}

\subsection*{T3 -- Kategorien}
\begin{enumerate}
  \item  \textbf{Ziel des Tasks (kurz, technisch).} Erweiterung des Todo-Datenmodells um Kategorisierung und Bereitstellung über API und UI, sodass Todos kategoriebasiert erstellt, gespeichert und angezeigt werden können.

  \item  \textbf{Typische Implementierungsvarianten (mit Häufigkeiten).} Ein separates Kategorie-Konzept (Model/Tabelle/Collection mit Relation via \texttt{categoryId}) dominiert (14/18). Inline-Kategorien als String sind selten (1/18). In 3/18 Branches ist das Kategorie-Konzept in Struktursignalen nicht klar erkennbar.

  \item  \textbf{Beobachtete Stabilitätsmuster (konsistent vs. inkonsistent).} Konsistent sind konsistente Persistenz und stabile Referenzen (IDs/Relation). Inkonsistent sind rein UI-seitige Kategorien oder inkonsistente Relation/Serialisierung.

  \item  \textbf{Beobachtete Implementationsmuster.} Häufig treten zusätzliche \texttt{Category}-Modelle und Routen auf. In einzelnen Branches werden Kategorien und Tags terminologisch vermischt; dabei variiert die Bezeichnung, während die Datenrepräsentation nicht in allen Fällen konsistent vereinheitlicht ist.

  \item  \textbf{Beobachtete Teilumsetzungen / offene Punkte.} In einzelnen Branches bleiben Kategorien rein UI-seitig, oder Persistenz/Serialisierung/Relation ist nicht durchgängig konsistent umgesetzt.
\end{enumerate}

\subsection*{T4 -- Dateianhänge (Attachments)}
\begin{enumerate}
  \item  \textbf{Ziel des Tasks (kurz, technisch).} Upload und Zuordnung von Dateien zu Todos (typischerweise \texttt{multipart/form-data}), Speicherung (Dateisystem oder DB) und Rückgabe verwertbarer Metadaten über die API.

  \item  \textbf{Typische Implementierungsvarianten (mit Häufigkeiten).} In 8/18 Branches ist ein Multer-/Upload-Ansatz erkennbar. In 10/18 Branches sind Attachments nicht oder nur indirekt oder unklar umgesetzt. Wo Upload existiert, ist Storage häufig lokal (Uploads-Ordner) oder metadatenbasiert (Dateiname/URL).

  \item  \textbf{Beobachtete Stabilitätsmuster (konsistent vs. inkonsistent).} Konsistent sind klare Upload-Verträge (Feldnamen/Metadaten) und stabile Zuordnung zum Todo. Inkonsistent sind uneinheitliche Response-Formate und fehlende Metadaten-Persistenz.

  \item  \textbf{Beobachtete Implementationsmuster.}\\
  Multer-Templates mit Boilerplate-Routen sind häufig. Daneben finden sich sowohl Varianten mit separaten Attachment-Entitäten als auch Upload-Varianten ohne stabilen Metadaten-Rückkanal.

  \item  \textbf{Beobachtete Teilumsetzungen / offene Punkte.} Wo Upload existiert, sind Response-Formate, Feldnamen und Metadatenpersistenz nicht in allen Branches konsistent.
\end{enumerate}

\subsection*{T5 -- Kalender / Multi-Day}
\begin{enumerate}
  \item  \textbf{Ziel des Tasks (kurz, technisch).} Abbildung von Todos im Kalender, inklusive Unterstützung von Zeiträumen (Multi-Day), sowie konsistente Query-/Filter-Semantik (z.B. Overlap-Logik) zur effizienten Darstellung in Kalenderansichten.

  \item  \textbf{Typische Implementierungsvarianten (mit Häufigkeiten).} In den Branches sind Range-Felder als Struktursignale verbreitet (z.B. \texttt{dueEndDate} neben \texttt{dueDate}). Gleichzeitig ist T5 in den (V3-)Integrationsläufen in 0/18 Branches als \textit{Execution observed} ausgewiesen.

  \item  \textbf{Beobachtete Stabilitätsmuster (konsistent vs. inkonsistent).} In den Artefakten treten zusätzliche Range-Felder häufig auf. Eine durchgängig erkennbare Range-Logik (Ende > Start) sowie eine konsistente Query-/Filter-Semantik (Overlap-Logik) ist in den Artefakten nicht in allen Branches erkennbar.

  \item  \textbf{Beobachtete Implementationsmuster.} Häufig treten UI-/Modell-Boilerplates (zusätzliche Datumsfelder, Kalender-Komponenten) sowie Range-Utilities auf. In den Artefakten ist der API-Vertrag (Parameter/Overlap-Regel) dabei nicht in allen Branches konsistent erkennbar.

  \item  \textbf{Beobachtete Teilumsetzungen / offene Punkte.} In den Artefakten sind Range-Felder und UI-Komponenten teilweise vorhanden, ohne dass eine durchgängig konsistente, durch \textit{Execution observed} ausgewiesene End-to-End-Semantik für Overlap/Filter und API-Vertrag vorliegt.
\end{enumerate}

\section{Code-Qualität und Metriken}
\label{sec:results_quality}

Dieses Unterkapitel präsentiert die objektiven Messwerte zur Code-Qualität, basierend auf automatisierten Analysetools.

% DATENQUELLEN:
% - ESLint Reports (Linting Errors)
% - Test Reports (Pass/Fail Raten)
% - Lighthouse/Playwright (falls vorhanden)
%
% HINWEISE:
% - Rein quantitative Darstellung
% - Keine Bewertung der Ursachen (das folgt in Kap. 7)

\section{Chat-Interaktionsmuster}
\label{sec:results_chat}

Die Analyse der Interaktionen zwischen Probanden und dem KI-Modell wird in diesem Abschnitt dargestellt.

% DATENQUELLEN:
% - Chat-Logs (Prompts, Antworten)
%
% HINWEISE:
% - Prompt-Qualität (Struktur, Kontext)
% - Debugging-Loops (Ping-Pong-Effekte)
% - Nutzung von Claude Sonnet 4.5 Spezifika

\section{Triangulation der Datenquellen}
\label{sec:results_triangulation}

Hier werden die Ergebnisse aus Code-Analyse, Chat-Logs und Test-Reports zusammengeführt, um Zusammenhänge sichtbar zu machen.

% DATENQUELLEN:
% - Aggregierte Daten aus allen vorherigen Analysen
%
% HINWEISE:
% - Zusammenhang zwischen Chat-Verhalten und Code-Qualität
% - Zusammenhang zwischen Strategie (z.B. DB-Wahl) und Testerfolg
% - Rein deskriptive Korrelationen, keine kausale Interpretation

\section{Zusammenfassung der Ergebnisse}
\label{sec:results_summary}

Kapitel~6 dokumentiert die Ergebnisse der Evaluation und stellt beobachtete Ausprägungen und Verteilungen dar, ohne diese zu bewerten oder zu erklären.
Die Ergebnisdarstellung basiert auf mehreren Datenquellen (Pipeline-Status \,\texttt{attempted}/\texttt{implemented}/\texttt{verified}, Code- und Test-Artefakte, Chat-Exports sowie Survey-Merkmale) und bleibt durchgängig deskriptiv.
Im Folgenden werden die zentralen Befundlinien systematisch zusammengeführt.

\paragraph{Pipeline-Ergebnisse entlang der Tasks (T1--T5).}
Über die Tasks hinweg zeigt sich ein gemeinsamer Trichter-Effekt: Attempted-Markierungen sind in der Breite vorhanden, während die Anteile der als Implemented$^\ast$ (Trichter-Evidenzstufe) bzw. Verified ausgewiesenen Ergebnisse mit zunehmendem Integrations- und Abhängigkeitsgrad abnehmen.
Dabei ist die Ausprägung dieses Trichters nicht in allen Branches gleichförmig, sondern zeigt Spannweiten zwischen sehr kurzem Task-Fortschritt und einem deutlich weiter reichenden Umsetzungsstand.
In der Gegenüberstellung der Aufgaben treten dabei zwei Bündel hervor:
\begin{itemize}
\item \textbf{Backend-nahe Tasks (T1/T2).} T1 ist in den Branches in unterschiedlichen Persistenzvarianten umgesetzt, wobei eine Variante dominiert und mehrere Alternativen vereinzelt auftreten. T2 wird in der Mehrzahl der Branches als Implemented$^\ast$ ausgewiesen und ist damit ein häufig realisiertes Backend-Element.
\item \textbf{Integrations- und Frontend-nahe Tasks (T3--T5).} Bei T3 und T4 nehmen die Implemented\textsubscript{static}-Ausweise ab. Zugleich werden bei einzelnen Tasks Diskrepanzen zwischen Implemented\textsubscript{static} (heuristische Artefakt-Markierung) und Verified sichtbar (ohne dass daraus in Kapitel~6 Schlussfolgerungen abgeleitet werden). T5 bleibt in den Artefakten durchgängig ohne Implemented\textsubscript{static}- und ohne Verified-Ausweise.
\end{itemize}
Über alle Tasks hinweg gilt: Verified, Implemented\textsubscript{static} und Attempted sind empirisch unterscheidbare Statuskategorien und fallen nicht notwendigerweise zusammen. Für die Trichterdarstellung wird Implemented$^\ast$ (Attempted $\wedge$ (Implemented\textsubscript{static} $\vee$ Verified)) als konsistente Evidenzstufe genutzt.
Kapitel~6 beschreibt diese Statusausweise als Ergebnisvariablen der Pipeline und nutzt sie als gemeinsame Referenz, um Code-, Test- und Chat-Artefakte pro Branch in Beziehung zu setzen.

\paragraph{Technische Qualität (aggregiert, ohne Kausalannahmen).}
Die in Kapitel~6 berichteten Qualitäts- und Technikindikatoren zeigen insgesamt eine heterogene Ausprägung über die Branches hinweg.
Dazu gehören (a) Unterschiede im Linting-Status und in der Häufigkeit typischer TypeScript-Surrogate (z.\,B. Any-Nutzung), (b) Spannweiten bei strukturbezogenen Metriken (z.\,B. Anzahl TypeScript-Dateien) sowie (c) variierende Testausgänge bis hin zu Konstellationen mit Fehlanteilen.
Diese Befunde werden in Kapitel~6 als Ko-Existenz von Pipeline-Fortschritt und technischen Merkmalen beschrieben. Es wird keine Korrelation oder Wirkbeziehung behauptet.
Zusätzlich wird sichtbar, dass sich technische Ausprägungen nicht auf einen einheitlichen "Projektzustand" reduzieren lassen: Branches mit ähnlichem Pipeline-Fortschritt können unterschiedliche Muster in Linting-, Typisierungs- und Testindikatoren aufweisen.
Zugleich werden die Messgrenzen betont: Die verwendeten Indikatoren ersetzen keine Coverage-, Komplexitäts- oder Performance-Metriken und erlauben entsprechend keine Aussagen zu nicht gemessenen Qualitätsdimensionen.

\paragraph{Chat-Interaktionsmuster (kompakt).}
Die Chat-Analysen zeigen große Spannweiten bei der Interaktionsintensität, insbesondere bei Wortanzahl und Prompt-Antwort-Zyklen.
Über die Branches hinweg koexistieren kurze, wenig strukturierte Prompts mit detaillierten Mehrpunkt-Prompts. Beide Formen treten im selben Datensatz nebeneinander auf.
Als wiederkehrendes Strukturmerkmal sind Debugging-Iterationen sichtbar (z.\,B. wiederholte Fehlermarker und erneute Anpassungsversuche), ohne dass daraus eine qualitative Einordnung ("gut/schlecht") abgeleitet wird.
Neben Intensitätsmaßen werden in Kapitel~6 auch Muster in der Interaktionsstruktur beschrieben, etwa wiederkehrende Task-Referenzen, Wechsel zwischen Aufgaben sowie Tooling-/Aktionsformulierungen im Verlauf der Chats.

\paragraph{Triangulation auf Meta-Ebene.}
Die Triangulation macht sichtbar, dass unterschiedliche Datenquellen teils kongruente, teils divergente Signale liefern: Pipeline-Status, technische Indikatoren, Testausgänge, Chat-Merkmale und Survey-Merkmale können für dieselben Branches zusammenpassen oder auseinanderlaufen.
Insbesondere verdeutlicht die Gegenüberstellung, dass Verified $\neq$ Implemented\textsubscript{static} $\neq$ Attempted und dass Spannungsfelder zwischen Statusausweisen, technischen Beobachtungen und Interaktionsmustern empirisch auftreten.
Auch Survey-Merkmale werden als kategoriale Kontextvariablen einbezogen (Selbsteinstufungsgruppen und benannte Hürdenkategorien), ohne dass daraus individuelle Zuschreibungen oder Bewertungen abgeleitet werden.
In der Gesamtschau entstehen damit nebeneinander stehende Ergebnislinien, die Gemeinsamkeiten und Unterschiede zwischen Branches sichtbar machen, ohne diese zu erklären.
Kapitel~6 stellt diese Spannungsfelder dar, löst sie jedoch nicht auf.

\paragraph{Abgrenzung zu Kapitel~7.}
Kapitel~6 beschreibt, \emph{was} beobachtet wurde. Kapitel~7 diskutiert, \emph{wie} diese Befunde einzuordnen sind und \emph{warum} bestimmte Muster auftreten könnten.

